% =====================================================================
%  CryptoVault — Phase 2 Completion Report
%  Compile free on https://www.overleaf.com (Upload -> Recompile)
%  or locally with MiKTeX:  pdflatex phase2_done.tex
% =====================================================================
\documentclass[11pt,a4paper]{article}

\usepackage[utf8]{inputenc}
\usepackage[T1]{fontenc}
\usepackage[margin=1in]{geometry}
\usepackage{lmodern}
\usepackage{microtype}
\usepackage{xcolor}
\usepackage{listings}
\usepackage{enumitem}
\usepackage{tcolorbox}
\usepackage{titlesec}
\usepackage{fancyhdr}
\usepackage[hidelinks,colorlinks=true,linkcolor=blue!60!black,urlcolor=teal]{hyperref}

\definecolor{accent}{RGB}{20,90,130}
\definecolor{okgreen}{RGB}{30,120,60}
\definecolor{lesson}{RGB}{95,55,140}
\definecolor{codebg}{RGB}{245,245,248}

\lstset{
  backgroundcolor=\color{codebg}, basicstyle=\ttfamily\small,
  breaklines=true, frame=single, rulecolor=\color{gray!40},
  framesep=6pt, xleftmargin=10pt, numbers=none, showstringspaces=false,
}

\newtcolorbox{done}{colback=okgreen!7,colframe=okgreen!55,boxrule=0.6pt,
  left=8pt,right=8pt,top=6pt,bottom=6pt,title=\textbf{Verified}}
\newtcolorbox{lessonbox}{colback=lesson!7,colframe=lesson!55,boxrule=0.6pt,
  left=8pt,right=8pt,top=6pt,bottom=6pt,title=\textbf{Lesson learned}}
\newtcolorbox{note}{colback=accent!6,colframe=accent!55,boxrule=0.6pt,
  left=8pt,right=8pt,top=6pt,bottom=6pt}

\pagestyle{fancy}\fancyhf{}
\fancyhead[L]{\small CryptoVault — Phase 2 Report}
\fancyhead[R]{\small\thepage}
\renewcommand{\headrulewidth}{0.4pt}

\titleformat{\section}{\large\bfseries\color{accent}}{\thesection}{0.6em}{}
\newcommand{\code}[1]{\texttt{#1}}

\title{\vspace{-1cm}\textbf{Phase 2 — Authentication Core}\\[3pt]
\large Completion Report}
\author{CryptoVault}
\date{21 June 2026}

\begin{document}
\maketitle
\thispagestyle{fancy}

\begin{note}
\textbf{Phase goal (from the plan):} register, login, and logout working with bcrypt-hashed
passwords, JWTs, and a Redis token blacklist --- and a logged-out token actually rejected on the
next request.\\[4pt]
\textbf{Status:} \textbf{Complete and verified live} against MySQL + Redis. 23 automated tests pass.
Shipped as \textbf{PR \#2} (\code{phase-2-auth} $\rightarrow$ \code{main}) at
\url{https://github.com/Jenak26/cryptovault/pull/2}. Builds on the teammate's merged Phase 1 schema.
\end{note}

% ---------------------------------------------------------------
\section{What "Phase 2" Means}
% ---------------------------------------------------------------
Phase 1 created the tables; Phase 2 is the first real \emph{feature}: letting a human prove who they
are. A user registers (their password is hashed, never stored in plain text), logs in (and receives
a signed token they attach to future requests), and logs out (their token is invalidated). The hard
part is doing it \emph{safely} --- the right error messages, no leaks, and tokens that can actually
be turned off.

% ---------------------------------------------------------------
\section{What Was Built}
% ---------------------------------------------------------------
\begin{itemize}[leftmargin=*]
  \item \textbf{Endpoints} --- \code{POST /api/auth/register}, \code{/login}, \code{/logout}, plus a
        protected \code{GET /api/me} to prove tokens are enforced.
  \item \textbf{\code{JwtService}} --- creates and validates JSON Web Tokens (library: \code{jjwt}
        0.12.6, algorithm HS256). Each token carries the user id, role, a unique id (\code{jti}), and
        an expiry. Lifetime: \textbf{1 hour}. The signing secret lives in \code{application.yml} but
        can be overridden by an environment variable for production.
  \item \textbf{\code{TokenBlacklist}} --- a Redis-backed list of logged-out token ids. On logout the
        token's \code{jti} is stored under \code{bl:jti:<id>} with a time-to-live equal to the
        token's remaining life, so it cleans itself up automatically.
  \item \textbf{\code{JwtAuthFilter}} --- runs on every request, reads the \code{Authorization:
        Bearer} header, validates the token, checks the blacklist, and tells Spring Security who the
        caller is (with authority \code{ROLE\_USER} / \code{ROLE\_ADMIN}).
  \item \textbf{\code{AuthService}} --- the brain: bcrypt hashing, the login check, and the logout
        blacklist call.
  \item \textbf{Supporting pieces} --- request/response DTOs (Java \code{record}s) with validation,
        a \code{GlobalExceptionHandler} mapping errors to clean JSON, a 401 entry point, and the
        \code{SecurityConfig} wiring.
  \item \textbf{\code{dev-console.html}} --- a small throwaway web page to click through the whole
        flow in a browser (dev-only; safe to delete before any real deploy).
\end{itemize}

% ---------------------------------------------------------------
\section{Key Decisions (locked)}
% ---------------------------------------------------------------
\begin{itemize}[leftmargin=*]
  \item \textbf{Generic login errors.} A wrong password \emph{and} an unknown email both return the
        same \code{401 "Invalid credentials"}. This stops an attacker from discovering which emails
        are registered. (Registration is the one exception --- a duplicate email must return
        \code{409}, because signup unavoidably reveals the address is taken.)
  \item \textbf{One short-lived token, no refresh tokens.} Refresh tokens are a Phase~V2 concern ---
        not built now (YAGNI).
  \item \textbf{Secret from the environment.} A dev default lets the app run out of the box; real
        deployments override \code{CRYPTOVAULT\_JWT\_SECRET}. Secrets never get committed.
  \item \textbf{Roles are uppercase} (\code{USER}/\code{ADMIN}) to match Phase 1's seed and Spring's
        \code{hasRole()} convention.
\end{itemize}

% ---------------------------------------------------------------
\section{Verification (evidence, not claims)}
% ---------------------------------------------------------------
\begin{done}
\textbf{1. Automated tests --- 23/23 green.} \code{JwtService} (5, incl. expired + tampered token
rejection), \code{TokenBlacklist} (3), \code{AuthService} (6), \code{AuthController} (4),
\code{JwtAuthFilter} (4), and the application context load (1).
\end{done}

\begin{done}
\textbf{2. The real flow, run live against MySQL + Redis:}
\begin{lstlisting}
register      -> 200 + JWT
login         -> 200 + JWT
GET /api/me   -> 200   (valid token accepted)
GET /api/me   -> 401   (no token)
logout        -> 200
GET /api/me   -> 401   (SAME token, now rejected)
wrong password-> 401   (generic message)
\end{lstlisting}
The database stored bcrypt hashes (\code{\$2a\$10\$...}, 60 chars) --- \textbf{zero plaintext
passwords}.
\end{done}

% ---------------------------------------------------------------
\section{What I Learned}
% ---------------------------------------------------------------

\begin{lessonbox}
\textbf{1. Passing unit tests do \emph{not} mean the feature works.}
Login passed every unit test, then failed the moment it ran for real with a
\code{LazyInitializationException}. The cause: the code read the user's \code{role} \emph{after} the
database session had closed (a JPA "lazy loading" trap, made stricter by our \code{open-in-view:
false} setting). The unit tests missed it because they used \emph{mock} objects, which are never
lazy. \\[4pt]
\textbf{Why it matters:} mocks prove your \emph{logic}; only a real database proves your
\emph{wiring}. The fix was one annotation (\code{@Transactional}, which keeps the session open), but
the lesson is bigger than the fix.
\end{lessonbox}

\begin{lessonbox}
\textbf{2. A phase's "done when" must be testable \emph{inside} that phase.}
The plan said Phase~2 was done when "logout $\rightarrow$ token rejected on next call" --- but the
piece that actually \emph{rejects} a token (the JWT filter) was scheduled for Phase~3, and there was
no protected route to be rejected from. The plan quietly contradicted itself. A blacklist that
nothing reads is untestable. \\[4pt]
\textbf{What we did:} pulled the basic filter and one protected endpoint (\code{/api/me}) into
Phase~2 so its own success criterion is real, and re-scoped Phase~3 to RBAC + rate limiting. Both
\code{docs/PLAN.md} and the design spec were updated with a dated note.
\end{lessonbox}

\begin{lessonbox}
\textbf{3. Security is mostly about what you \emph{don't} reveal.}
The most important line of the login code is the one that makes a wrong password and a missing
account look identical. Good security here is a deliberate \emph{absence} of information --- no
existence leaks, no secrets in logs, no plaintext in the database.
\end{lessonbox}

% ---------------------------------------------------------------
\section{How to Do This Better Next Time}
% ---------------------------------------------------------------
\begin{enumerate}[leftmargin=*]
  \item \textbf{Run one real end-to-end test \emph{early}}, not just at the end. Had I hit the live
        endpoint after the first login implementation, the lazy-loading bug would have surfaced in
        minutes instead of after the whole service was written.
  \item \textbf{Before starting a phase, read its "done when" and ask: can I prove this with only
        what's in this phase?} If proving it needs something from a later phase, the plan is wrong ---
        fix the plan first, then code.
  \item \textbf{Keep the brainstorm $\rightarrow$ spec $\rightarrow$ plan $\rightarrow$ build order.}
        Writing the design spec first (what, why, which errors, which status codes) made the actual
        coding fast and removed guesswork. The spec lives in
        \code{docs/superpowers/specs/}.
  \item \textbf{Commit in small, labelled steps} (design, then core, then enforcement). When
        something breaks, a small commit is easy to inspect and revert.
  \item \textbf{Treat "it compiles" and "tests pass" as necessary but not sufficient.} The bar is a
        live run with the real database and cache.
\end{enumerate}

% ---------------------------------------------------------------
\section{Next: Phase 3 --- RBAC \& Rate Limiting}
% ---------------------------------------------------------------
With authentication proven, Phase~3 adds \emph{authorization}: role-based rules
(\code{@PreAuthorize("hasRole('ADMIN')")}) so admin-only routes return 403 for normal users, plus a
Redis-based rate limiter that blocks brute-force login attempts. The JWT filter it needs is already
built and shipping in this phase, so Phase~3 can focus purely on the role rules and the limiter.

\vspace{1em}
\begin{center}\itshape Phase 2 complete --- users can register, log in, and log out, and a
logged-out token is genuinely rejected. Verified end-to-end, not just asserted.\end{center}

\end{document}
